\chapter{第三讲\quad 犯罪客体}
\label{lecture:03}
\noindent{\small\hyperref[ref:lecture:03]{查阅本讲扩展规则}}\par

\section{犯罪客体}

\paragraph{定义\quad 行为对象与法益分属两个层次}
设 \(\mathcal O\) 为行为可指向的对象集合，\(\mathcal B\) 为刑法保护利益的类型集合。用 \(\mathrm{Target}(x,o)\) 表示事件 \(x\) 作用于对象 \(o\in\mathcal O\)，用 \(\mathrm{Carries}(o,b)\) 表示对象 \(o\) 承载法益 \(b\in\mathcal B\)。一个对象可承载多种利益，同一种利益也可有许多载体。再定义 \(\mathrm{Protects}(g,b)\) 表示罪名 \(g\) 保护法益 \(b\)，\(\mathrm{Hits}(x,b)\) 表示事件侵害或者在法定意义上危及 \(b\)。教材的法益分析可写成
\[
 \mathrm{Relevant}(g,x)
 \;:=\;\exists b\,
    \bigl(\mathrm{Protects}(g,b)\land\mathrm{Hits}(x,b)\bigr).
\]
这里 \(\mathrm{Relevant}\) 是分析方向，罪名成立仍需核实具体构成。第2条、第13条列出刑法保护任务及犯罪概念的法定边界。以法益侵害说明犯罪的实质是教材的理论立场。

\paragraph{推导\quad 手机相同，侵害方式不同}
甲秘密将乙的手机从乙的控制下取走并据为己有。以 \(o\) 为手机，以 \(b\) 为乙对财物受刑法保护的利益。由所有与控制事实得 \(\mathrm{Carries}(o,b)\)，由取走行为和失控事实得 \(\mathrm{Hits}(x,b)\)。第264条的保护目的与 \(b\) 相符，可继续考察窃取、主观目的与门槛。若仅改为把手机砸坏并留在原处，则 \(o,b\) 可仍相同，而取走、占有转移与非法占有目的发生变化。破坏财物案应考察第275条的毁坏方式与其余条件。这个对照给出反例：相同的对象和法益不能推出相同的罪名。

\paragraph{争议推导\quad 住宅的保护目的}
以第245条的非法侵入住宅分支为例。题设一处房屋长期无人生活，但产权人保有可现实行使的进入许可权。行为人未经许可进入并滞留。教材列出两个可比较的保护目的。以现实居住安宁为中心的解释把生活安宁是否受扰视为关键；若题设排除现实居住安宁，按此立场，该要件不能成立。以进入许可权为中心的解释重视许可权受侵害；在该房屋依此立场仍属受保护住宅的前提下，无权进入与滞留可支持侵害该法益。推导可列为
\begin{align*}
  F\land\mathrm{Focus}_{\text{安宁}}\land\neg\mathrm{Living}
   &\Rightarrow \neg\mathrm{Hits}(x,\text{居住安宁}),\\
  F\land\mathrm{Focus}_{\text{许可}}\land\mathrm{Residence}
   &\Rightarrow \mathrm{Hits}(x,\text{进入许可利益}).
\end{align*}
\(F\) 是共同事实，\(\mathrm{Living}\) 指题设存在现实居住生活，\(\mathrm{Residence}\) 指依第二种立场确认为住宅。第二条只推进法益分析，继续得出有罪仍需逐项审查构成及证明。若改为住户仅当天上班，便无法从进门时「无人」推出长期没有居住安宁。这是事实推导的边界。

\section{保障人权}

\paragraph{法定条件先于保护目的}
令 \(K_g(x)\) 表示已经证立的法定入罪条件，\(P_g(x)\) 表示处罚对保护法益有利。即使承认 \(P_g(x)\)，也必须经过 \(K_g(x)\) 的检验。第3条的约束可写成
\[
 \neg K_g(x)\ \Rightarrow\ \text{依现有条件不能以 }g\text{ 定罪}.
\]
此处 \(\neg K_g(x)\) 指该条件合取没有成立，需在答案中进一步说明是法律上已证实不符合，还是事实尚未证明。第13条对显著轻微、危害不大的情形作出限制，应用时须核实具体情节。

\paragraph{推导\quad 脱逃题先看关押条件}
自拟两案，均有人脱离现场控制。第一案甲是依法被关押的犯罪嫌疑人，故意摆脱监管并逃离，题设其余要件均满足。第316条要求的主体身份、依法关押与脱逃事实均可对应，进入脱逃罪证明。第二案乙被私人非法拘禁，逃离行为虽扰乱拘禁者控制，但「依法被关押」一项明确不成立。第316条的合取因此为假。若实际题设属于曾被依法关押、后来发现原案错判的案件，须另核原关押的法律基础、当时的救济与其他出罪理由，不能仅凭后来改判就套用私人非法拘禁案。

\subsection*{闭卷复述与检验}
\textbf{口诀}\quad「对象看东西，客体看利益。先法定，后归责。」复述时给一个案情画三列，依次填对象、法益、法定动词。遇到目的解释争议，在末尾写清改变的是哪一个要件的解释。

\begin{enumerate}
\item 手机被拿走和被砸坏，哪个量可保持相同，哪个必须重算？\\
\textbf{答}\quad 手机及其所承载的财产法益可保持相同。需重新判断窃取、毁坏、占有变化、故意内容及各罪的法定门槛。
\item 某人破坏了他人对自己的非法控制。为了维护秩序，能否用脱逃罪处罚？\\
\textbf{答}\quad 第316条首先要求依法被关押的法定主体。私人非法拘禁不满足该项，保护目的不能补齐它。
\item 住户两天出差，能否由门内暂时无人，推出没有住宅生活安宁？\\
\textbf{答}\quad 不能。暂时离开与长期未实际用于居住是不同事实，须据房屋用途、居住连续性与许可状态重新分析。
\end{enumerate}
